\clearpage

% APENDICE C
\section{Áreas de Aprendizado e Subtópicos}

\subsection*{Área 1 --- Lógica e Raciocínio Matemático}

\subsubsection*{Proposições e Conectivos Lógicos}
Proposições são afirmações que podem ser verdadeiras ou falsas. Os conectivos lógicos
(negação $\neg$, conjunção $\wedge$, disjunção $\vee$, implicação $\Rightarrow$ e
bicondicional $\Leftrightarrow$) permitem combinar proposições para formar argumentos
complexos.

\subsubsection*{Tabelas-Verdade e Equivalências}
Tabelas-verdade analisam o comportamento de expressões lógicas em todos os casos possíveis.
As Leis de De Morgan ($\neg(A \wedge B) \equiv \neg A \vee \neg B$) e outras equivalências
permitem simplificar expressões.

\subsubsection*{Quantificadores e Predicados}
Quantificadores universais ($\forall$, ``para todo'') e existenciais ($\exists$, ``existe'')
permitem expressar afirmações sobre conjuntos de objetos. A lógica de predicados forma a
base da linguagem matemática formal.

\subsubsection*{Técnicas de Demonstração}
Técnicas como demonstração direta, por contra positiva, por contradição e por indução são
ferramentas fundamentais do raciocínio matemático.

\subsection*{Área 2 --- Matemática Discreta}

\subsubsection*{Teoria Ingênua dos Conjuntos}
Operações como união ($A \cup B$), interseção ($A \cap B$), diferença e produto cartesiano
($A \times B$) aparecem constantemente em estruturas de dados e teoria de tipos.

\subsubsection*{Combinatória}
Permutações $P(n,k)$, combinações $\binom{n}{k}$ e o princípio da inclusão-exclusão
aparecem em análise de complexidade de algoritmos e contagem de estruturas.

\subsubsection*{Relações e Funções}
Funções são relações com propriedades específicas essenciais para compreender mapeamentos
em programação: injetividade (saídas distintas para entradas distintas), sobrejetividade
(todo elemento do contradomínio é atingido) e bijetividade (ambas as propriedades).

\subsubsection*{Teoria dos Grafos}
Grafos $G = (V, E)$ modelam relações entre objetos. Conceitos como caminhos, ciclos,
conectividade e árvores têm aplicações diretas em redes e algoritmos de busca.

\subsubsection*{Indução Matemática}
Técnica de demonstração para afirmações sobre números naturais, especialmente útil para
provar a correção de algoritmos recursivos.

\subsection*{Área 3 --- Álgebra e Aritmética}

\subsubsection*{Aritmética Modular}
$a \equiv b \pmod{m}$ estuda o comportamento dos números em relação a um módulo.
Fundamental em criptografia e hashing.

\subsubsection*{Álgebra Linear Básica}
Vetores, matrizes e transformações lineares são essenciais para computação gráfica,
aprendizado de máquina e física computacional.

\subsubsection*{Teoria dos Números}
Divisibilidade, números primos, $\gcd(a,b)$, $\text{lcm}(a,b)$ e o Algoritmo de Euclides
são a base de algoritmos criptográficos como RSA.

\subsection*{Área 4 --- Probabilidade e Estatística}

\subsubsection*{Probabilidade Básica}
$P(A) \in [0,1]$. Probabilidade condicional $P(A|B) = \frac{P(A \cap B)}{P(B)}$ e
independência são fundamentais para algoritmos aleatorizados.

\subsubsection*{Estatística Descritiva}
Médias, medianas, desvios padrão $\sigma$, percentis e visualizações são indispensáveis
para qualquer projeto de análise de dados.

\subsubsection*{Variáveis Aleatórias e Distribuições}
Distribuições binomial, Poisson e normal descrevem padrões comuns em dados reais. Esperança
$E[X]$ e variância $\text{Var}(X)$ são medidas centrais.

\subsection*{Área 5 --- Algoritmos e Complexidade}

\subsubsection*{Recorrências}
O Teorema Mestre resolve recorrências da forma $T(n) = aT(n/b) + f(n)$, permitindo analisar
algoritmos recursivos como Merge Sort
($T(n) = 2T(n/2) + \mathcal{O}(n) = \mathcal{O}(n \log n)$).

\subsubsection*{Análise de Algoritmos}
A notação Big-$\mathcal{O}$ permite comparar algoritmos de forma independente de hardware:
$\mathcal{O}(1) \subset \mathcal{O}(\log n) \subset \mathcal{O}(n) \subset
\mathcal{O}(n \log n) \subset \mathcal{O}(n^2) \subset \mathcal{O}(2^n)$.

\subsubsection*{Técnicas de Design de Algoritmos}
Divisão e conquista, programação dinâmica e algoritmos gulosos são paradigmas fundamentais
para resolver problemas computacionais eficientemente.
